% --- Archon physics formalization source begin ---
% archon:physics
% archon:covers PhyXMiniProblems/problem_phyx_mini_0241.lean
% archon:source-report reports/phyx_mini/problem_phyx_mini_0241.source.json

\chapter{Physics problem phyx\_mini\_0241}
\label{ch:PhyXMiniProblems_problem_phyx_mini_0241}

\paragraph{Problem source.}
\#\# Physical scenario

Tension is maintained in a string as in figure. The observed wave speed is \textbackslash{}( v = 24.0\textbackslash{},\textbackslash{}mathrm\{m/s\} \textbackslash{}) when the suspended mass is \textbackslash{}( m = 3.00\textbackslash{},\textbackslash{}mathrm\{kg\} \textbackslash{}).The suspended mass is \textbackslash{}( m = 2.00\textbackslash{},\textbackslash{}mathrm\{kg\} \textbackslash{}).

\#\# Question

What is the wave speed?

\#\# Answer choices

- A: A: 22.6m/s

- B: B: 21.6m/s

- C: C: 20.6m/s

- D: D: 19.6m/s

\#\# Figure caption (auxiliary; use the image as primary evidence)

The image depicts a physics-related setup involving a pulley system. There is a table with a pulley attached to its edge. A rope runs over the pulley, one end tied to a fixed point on the left and the other to a square object labeled "m" hanging off the table. The setup illustrates basic principles of mechanics, such as tension and gravitational force.

\#\# Dataset metadata

Wave Properties; Multi-Formula Reasoning, Physical Model Grounding Reasoning

\paragraph{Recorded answer/context.}
D: D: 19.6m/s

\paragraph{Figure/image path.}
/root/proposal\_for\_physic/science-mango/phyx\_mini\_run/phyx\_data/test\_image/241.png

\paragraph{Formalization target.}
create a compiling Lean file with sorry bodies at `PhyXMiniProblems/problem\_phyx\_mini\_0241.lean`. The Lean declarations must preserve the physical quantities, dimensions or dimensional roles, figure labels, governing-law hypotheses, and final relation expressed by this problem.
Use Mathlib/Physlib names found through LeanExplore where available. If a physics API is missing, introduce faithful local abstractions rather than scalar placeholder aliases.

\begin{theorem}[Physics formalization target]
\label{thm:physics:phyx_mini_0241:target}
\lean{PhyXMiniProblems.ProblemPhyXMini0241.problem_phyx_mini_0241}
\uses{def:physics:phyx-mini-0241:phyxminiproblems-problemphyxmini0241-stringpulleysetup, def:physics:phyx-mini-0241:phyxminiproblems-problemphyxmini0241-matchesproblemandprimaryfigure, def:physics:phyx-mini-0241:phyxminiproblems-problemphyxmini0241-haspositivephysicalparameters, def:physics:phyx-mini-0241:phyxminiproblems-problemphyxmini0241-satisfiestautstringandidealpulleylaws, def:physics:phyx-mini-0241:phyxminiproblems-problemphyxmini0241-recordedanswerchoice, def:physics:phyx-mini-0241:phyxminiproblems-problemphyxmini0241-matchesanswerchoice}
The assigned autoformalize agent should translate this physics problem into Lean declarations in the covered file, with theorem and lemma proof bodies written as `by sorry`.
\end{theorem}
\begin{proof}
This is an autoformalization task, not a proof task. Produce statements that can later be proved by the physics prover without weakening the physical model.
\end{proof}

% --- Archon declaration topology begin ---
\section{Declaration topology}
\begin{definition}[Mass Quantity]
  \label{def:physics:phyx-mini-0241:phyxminiproblems-problemphyxmini0241-massquantity}
  \lean{PhyXMiniProblems.ProblemPhyXMini0241.MassQuantity}
  A nonnegative physical mass
\end{definition}
\begin{proof}
  This is the declared model definition of Mass Quantity.
\end{proof}
\begin{definition}[Acceleration Quantity]
  \label{def:physics:phyx-mini-0241:phyxminiproblems-problemphyxmini0241-accelerationquantity}
  \lean{PhyXMiniProblems.ProblemPhyXMini0241.AccelerationQuantity}
  A nonnegative physical acceleration magnitude
\end{definition}
\begin{proof}
  This is the declared model definition of Acceleration Quantity.
\end{proof}
\begin{definition}[Tension Quantity]
  \label{def:physics:phyx-mini-0241:phyxminiproblems-problemphyxmini0241-tensionquantity}
  \lean{PhyXMiniProblems.ProblemPhyXMini0241.TensionQuantity}
  A nonnegative force, used here for string tension
\end{definition}
\begin{proof}
  This is the declared model definition of Tension Quantity.
\end{proof}
\begin{definition}[Linear Mass Density Quantity]
  \label{def:physics:phyx-mini-0241:phyxminiproblems-problemphyxmini0241-linearmassdensityquantity}
  \lean{PhyXMiniProblems.ProblemPhyXMini0241.LinearMassDensityQuantity}
  A nonnegative mass per unit length of string
\end{definition}
\begin{proof}
  This is the declared model definition of Linear Mass Density Quantity.
\end{proof}
\begin{definition}[Speed Quantity]
  \label{def:physics:phyx-mini-0241:phyxminiproblems-problemphyxmini0241-speedquantity}
  \lean{PhyXMiniProblems.ProblemPhyXMini0241.SpeedQuantity}
  A unit-independent propagation speed
\end{definition}
\begin{proof}
  This is the declared model definition of Speed Quantity.
\end{proof}
\begin{definition}[quantity Readout]
  \label{def:physics:phyx-mini-0241:phyxminiproblems-problemphyxmini0241-quantityreadout}
  \lean{PhyXMiniProblems.ProblemPhyXMini0241.quantityReadout}
  Read a nonnegative dimensionful quantity in a coherent unit system
\end{definition}
\begin{proof}
  This is the declared model definition of quantity Readout.
\end{proof}
\begin{definition}[mass In Kilograms]
  \label{def:physics:phyx-mini-0241:phyxminiproblems-problemphyxmini0241-massinkilograms}
  \lean{PhyXMiniProblems.ProblemPhyXMini0241.massInKilograms}
  \uses{def:physics:phyx-mini-0241:phyxminiproblems-problemphyxmini0241-massquantity, def:physics:phyx-mini-0241:phyxminiproblems-problemphyxmini0241-quantityreadout}
  Kilogram readout of a physical mass
\end{definition}
\begin{proof}
  This is the declared model definition of mass In Kilograms.
\end{proof}
\begin{definition}[acceleration In Meters Per Second Squared]
  \label{def:physics:phyx-mini-0241:phyxminiproblems-problemphyxmini0241-accelerationinmeterspersecondsquared}
  \lean{PhyXMiniProblems.ProblemPhyXMini0241.accelerationInMetersPerSecondSquared}
  \uses{def:physics:phyx-mini-0241:phyxminiproblems-problemphyxmini0241-accelerationquantity, def:physics:phyx-mini-0241:phyxminiproblems-problemphyxmini0241-quantityreadout}
  Metres-per-second-squared readout of an acceleration
\end{definition}
\begin{proof}
  This is the declared model definition of acceleration In Meters Per Second Squared.
\end{proof}
\begin{definition}[tension In Newtons]
  \label{def:physics:phyx-mini-0241:phyxminiproblems-problemphyxmini0241-tensioninnewtons}
  \lean{PhyXMiniProblems.ProblemPhyXMini0241.tensionInNewtons}
  \uses{def:physics:phyx-mini-0241:phyxminiproblems-problemphyxmini0241-tensionquantity, def:physics:phyx-mini-0241:phyxminiproblems-problemphyxmini0241-quantityreadout}
  Newton readout of a string tension
\end{definition}
\begin{proof}
  This is the declared model definition of tension In Newtons.
\end{proof}
\begin{definition}[linear Mass Density In Kilograms Per Meter]
  \label{def:physics:phyx-mini-0241:phyxminiproblems-problemphyxmini0241-linearmassdensityinkilogramspermeter}
  \lean{PhyXMiniProblems.ProblemPhyXMini0241.linearMassDensityInKilogramsPerMeter}
  \uses{def:physics:phyx-mini-0241:phyxminiproblems-problemphyxmini0241-linearmassdensityquantity, def:physics:phyx-mini-0241:phyxminiproblems-problemphyxmini0241-quantityreadout}
  Kilograms-per-metre readout of a string's linear mass density
\end{definition}
\begin{proof}
  This is the declared model definition of linear Mass Density In Kilograms Per Meter.
\end{proof}
\begin{definition}[speed In Meters Per Second]
  \label{def:physics:phyx-mini-0241:phyxminiproblems-problemphyxmini0241-speedinmeterspersecond}
  \lean{PhyXMiniProblems.ProblemPhyXMini0241.speedInMetersPerSecond}
  \uses{def:physics:phyx-mini-0241:phyxminiproblems-problemphyxmini0241-speedquantity, def:physics:phyx-mini-0241:phyxminiproblems-problemphyxmini0241-quantityreadout}
  Metres-per-second readout of a wave speed
\end{definition}
\begin{proof}
  This is the declared model definition of speed In Meters Per Second.
\end{proof}
\begin{definition}[Trial]
  \label{def:physics:phyx-mini-0241:phyxminiproblems-problemphyxmini0241-trial}
  \lean{PhyXMiniProblems.ProblemPhyXMini0241.Trial}
  The measured calibration state and the state asked about in the problem
\end{definition}
\begin{proof}
  This is the declared model definition of Trial.
\end{proof}
\begin{definition}[Figure Component]
  \label{def:physics:phyx-mini-0241:phyxminiproblems-problemphyxmini0241-figurecomponent}
  \lean{PhyXMiniProblems.ProblemPhyXMini0241.FigureComponent}
  Individually identifiable components in the supplied figure
\end{definition}
\begin{proof}
  This is the declared model definition of Figure Component.
\end{proof}
\begin{definition}[String Route]
  \label{def:physics:phyx-mini-0241:phyxminiproblems-problemphyxmini0241-stringroute}
  \lean{PhyXMiniProblems.ProblemPhyXMini0241.StringRoute}
  The route followed by the string in the primary figure
\end{definition}
\begin{proof}
  This is the declared model definition of String Route.
\end{proof}
\begin{definition}[String Endpoint]
  \label{def:physics:phyx-mini-0241:phyxminiproblems-problemphyxmini0241-stringendpoint}
  \lean{PhyXMiniProblems.ProblemPhyXMini0241.StringEndpoint}
  The two ends of the pictured string and their attachments
\end{definition}
\begin{proof}
  This is the declared model definition of String Endpoint.
\end{proof}
\begin{definition}[Endpoint Attachment]
  \label{def:physics:phyx-mini-0241:phyxminiproblems-problemphyxmini0241-endpointattachment}
  \lean{PhyXMiniProblems.ProblemPhyXMini0241.EndpointAttachment}
  Objects attached to the two ends of the string
\end{definition}
\begin{proof}
  This is the declared model definition of Endpoint Attachment.
\end{proof}
\begin{definition}[String Pulley Setup]
  \label{def:physics:phyx-mini-0241:phyxminiproblems-problemphyxmini0241-stringpulleysetup}
  \lean{PhyXMiniProblems.ProblemPhyXMini0241.StringPulleySetup}
  \uses{def:physics:phyx-mini-0241:phyxminiproblems-problemphyxmini0241-massquantity, def:physics:phyx-mini-0241:phyxminiproblems-problemphyxmini0241-accelerationquantity, def:physics:phyx-mini-0241:phyxminiproblems-problemphyxmini0241-tensionquantity, def:physics:phyx-mini-0241:phyxminiproblems-problemphyxmini0241-linearmassdensityquantity, def:physics:phyx-mini-0241:phyxminiproblems-problemphyxmini0241-speedquantity, def:physics:phyx-mini-0241:phyxminiproblems-problemphyxmini0241-trial, def:physics:phyx-mini-0241:phyxminiproblems-problemphyxmini0241-figurecomponent, def:physics:phyx-mini-0241:phyxminiproblems-problemphyxmini0241-stringroute, def:physics:phyx-mini-0241:phyxminiproblems-problemphyxmini0241-stringendpoint, def:physics:phyx-mini-0241:phyxminiproblems-problemphyxmini0241-endpointattachment}
  Independent physical quantities and categorical figure data for the two trials. The single stringLinearMassDensity and gravitationalAcceleration fields express that the same string and location are used in both trials. Neither requested-speed value is assigned in this structure
\end{definition}
\begin{proof}
  This is the declared model definition of String Pulley Setup.
\end{proof}
\begin{definition}[Matches Problem And Primary Figure]
  \label{def:physics:phyx-mini-0241:phyxminiproblems-problemphyxmini0241-matchesproblemandprimaryfigure}
  \lean{PhyXMiniProblems.ProblemPhyXMini0241.MatchesProblemAndPrimaryFigure}
  \uses{def:physics:phyx-mini-0241:phyxminiproblems-problemphyxmini0241-massinkilograms, def:physics:phyx-mini-0241:phyxminiproblems-problemphyxmini0241-speedinmeterspersecond, def:physics:phyx-mini-0241:phyxminiproblems-problemphyxmini0241-stringpulleysetup}
  The categorical information read from the primary figure and the numerical data stated in the problem: 3.00 kg with measured speed 24.0 m/s, followed by a 2.00 kg load. In particular, the requested trial's speed is absent
\end{definition}
\begin{proof}
  This is the declared model definition of Matches Problem And Primary Figure.
\end{proof}
\begin{definition}[Has Positive Physical Parameters]
  \label{def:physics:phyx-mini-0241:phyxminiproblems-problemphyxmini0241-haspositivephysicalparameters}
  \lean{PhyXMiniProblems.ProblemPhyXMini0241.HasPositivePhysicalParameters}
  \uses{def:physics:phyx-mini-0241:phyxminiproblems-problemphyxmini0241-massinkilograms, def:physics:phyx-mini-0241:phyxminiproblems-problemphyxmini0241-accelerationinmeterspersecondsquared, def:physics:phyx-mini-0241:phyxminiproblems-problemphyxmini0241-tensioninnewtons, def:physics:phyx-mini-0241:phyxminiproblems-problemphyxmini0241-linearmassdensityinkilogramspermeter, def:physics:phyx-mini-0241:phyxminiproblems-problemphyxmini0241-speedinmeterspersecond, def:physics:phyx-mini-0241:phyxminiproblems-problemphyxmini0241-stringpulleysetup}
  Positivity and nondegeneracy conditions for the taut-string model
\end{definition}
\begin{proof}
  This is the declared model definition of Has Positive Physical Parameters.
\end{proof}
\begin{definition}[Satisfies Taut String And Ideal Pulley Laws]
  \label{def:physics:phyx-mini-0241:phyxminiproblems-problemphyxmini0241-satisfiestautstringandidealpulleylaws}
  \lean{PhyXMiniProblems.ProblemPhyXMini0241.SatisfiesTautStringAndIdealPulleyLaws}
  \uses{def:physics:phyx-mini-0241:phyxminiproblems-problemphyxmini0241-massinkilograms, def:physics:phyx-mini-0241:phyxminiproblems-problemphyxmini0241-accelerationinmeterspersecondsquared, def:physics:phyx-mini-0241:phyxminiproblems-problemphyxmini0241-tensioninnewtons, def:physics:phyx-mini-0241:phyxminiproblems-problemphyxmini0241-linearmassdensityinkilogramspermeter, def:physics:phyx-mini-0241:phyxminiproblems-problemphyxmini0241-speedinmeterspersecond, def:physics:phyx-mini-0241:phyxminiproblems-problemphyxmini0241-stringpulleysetup}
  The standard idealizations used for each trial: a frictionless, massless pulley transmits the same tension to both segments; static balance makes the vertical tension equal to the hanging weight; a uniform taut string obeys μ v² = T. These are general laws relating independent quantities. They contain neither the requested exact speed nor the displayed answer 19.6 m/s
\end{definition}
\begin{proof}
  This is the declared model definition of Satisfies Taut String And Ideal Pulley Laws.
\end{proof}
\begin{lemma}[requested Wave Speed exact]
  \label{lem:physics:phyx-mini-0241:phyxminiproblems-problemphyxmini0241-requestedwavespeed-exact}
  \lean{PhyXMiniProblems.ProblemPhyXMini0241.requestedWaveSpeed_exact}
  \uses{def:physics:phyx-mini-0241:phyxminiproblems-problemphyxmini0241-speedinmeterspersecond, def:physics:phyx-mini-0241:phyxminiproblems-problemphyxmini0241-stringpulleysetup, def:physics:phyx-mini-0241:phyxminiproblems-problemphyxmini0241-matchesproblemandprimaryfigure, def:physics:phyx-mini-0241:phyxminiproblems-problemphyxmini0241-haspositivephysicalparameters, def:physics:phyx-mini-0241:phyxminiproblems-problemphyxmini0241-satisfiestautstringandidealpulleylaws}
  Because the two trials use the same string and gravitational field, the laws give v\_requested / v\_calibration = sqrt (2 / 3). With the measured calibration speed, the requested unrounded value is 24 * sqrt (2 / 3) m/s
\end{lemma}
\begin{proof}
  The conclusion follows from the governing relations and figure readouts named in the statement of requested Wave Speed exact.
\end{proof}
\begin{definition}[Answer Choice]
  \label{def:physics:phyx-mini-0241:phyxminiproblems-problemphyxmini0241-answerchoice}
  \lean{PhyXMiniProblems.ProblemPhyXMini0241.AnswerChoice}
  Labels of the four answers printed in the problem
\end{definition}
\begin{proof}
  This is the declared model definition of Answer Choice.
\end{proof}
\begin{definition}[speed In Meters Per Second]
  \label{def:physics:phyx-mini-0241:phyxminiproblems-problemphyxmini0241-answerchoice-speedinmeterspersecond}
  \lean{PhyXMiniProblems.ProblemPhyXMini0241.AnswerChoice.speedInMetersPerSecond}
  \uses{def:physics:phyx-mini-0241:phyxminiproblems-problemphyxmini0241-speedinmeterspersecond, def:physics:phyx-mini-0241:phyxminiproblems-problemphyxmini0241-answerchoice}
  Wave speed printed beside each answer label, in metres per second
\end{definition}
\begin{proof}
  This is the declared model definition of speed In Meters Per Second.
\end{proof}
\begin{definition}[recorded Answer Choice]
  \label{def:physics:phyx-mini-0241:phyxminiproblems-problemphyxmini0241-recordedanswerchoice}
  \lean{PhyXMiniProblems.ProblemPhyXMini0241.recordedAnswerChoice}
  \uses{def:physics:phyx-mini-0241:phyxminiproblems-problemphyxmini0241-answerchoice}
  The answer label recorded in the supplied dataset metadata
\end{definition}
\begin{proof}
  This is the declared model definition of recorded Answer Choice.
\end{proof}
\begin{definition}[Matches Answer Choice]
  \label{def:physics:phyx-mini-0241:phyxminiproblems-problemphyxmini0241-matchesanswerchoice}
  \lean{PhyXMiniProblems.ProblemPhyXMini0241.MatchesAnswerChoice}
  \uses{def:physics:phyx-mini-0241:phyxminiproblems-problemphyxmini0241-speedinmeterspersecond, def:physics:phyx-mini-0241:phyxminiproblems-problemphyxmini0241-stringpulleysetup, def:physics:phyx-mini-0241:phyxminiproblems-problemphyxmini0241-answerchoice}
  Agreement with a speed displayed to one decimal place. The half-unit in the last printed digit is 0.05 m/s; using a tolerance avoids falsely asserting that the unrounded physical speed is exactly the displayed decimal
\end{definition}
\begin{proof}
  This is the declared model definition of Matches Answer Choice.
\end{proof}
% --- Archon declaration topology end ---
% --- Archon physics formalization source end ---
